\iflongversion
\subsection{What produces the compiled channel's hold}
\label{sec:isolation}
\phantomsection\label{sec:isolation-steering}\phantomsection\label{sec:isolation-frozen}%
\phantomsection\label{sec:isolation-lora}\phantomsection\label{sec:isolation-further}%
\phantomsection\label{sec:isolation-rank}\phantomsection\label{sec:isolation-amplitude}%
\phantomsection\label{sec:isolation-mmlu}

The compiled channel's hold comes from a fixed per-belief update, not from the modulation
responding to its input, and what makes that update effective is how it was trained.
Input-responsiveness is not the mechanism: frozen to a constant per-layer offset the modulation
still holds 88.9\% [81.1, 95.3] against an 86.9\% anchor. Neither is rank: truncating the
trained rank-16 update to rank 1 leaves persistence flat. Nor the context slot: modulation with
the slot empty holds 91.9\% with zero belief tokens in context, while the slot alone holds
57.5\%, matching context because it is context. What remains is close to a fixed per-belief
low-rank update, which is what a LoRA parameterizes; the LoRA controls that test this
are in Section~\ref{sec:objections} and Appendix~\ref{app:loramargin}.

What separates them after that is the training: the write function is trained across many
beliefs to prefer the assigned side in the presence of opposing content, so it resists opposing
content. A parameter- and objective-matched adapter recovers part of the persistence and not
the rest, and the remainder is unattributed between architecture and curriculum. On that reading the compiled channel's
contribution is a training recipe whose product is deliverable as a fixed low-rank delta, and
it remains the case the instrument is most useful on: a method that raises hold about as well
as anything tested and supplies no judgment with it. Steering from the same content is reported
as degenerate rather than as a baseline (14\% against the unconditioned host's 47\%), and
general knowledge is unchanged (MMLU 76.5\% vs.\ 74.8\%, McNemar \(p=0.11\)). Protocols and
per-ablation numbers are in Appendices~\ref{app:dstar} and~\ref{app:loramargin}.
\fi
